% Original PDF: Section 1, pages 1--2. Citations retain their printed text.
\section{Introduction}
Policy iteration alternates between evaluating a policy and improving it using the resulting
value estimates. In its optimistic form, the policy is improved before evaluation has converged.
Monte Carlo Exploring Starts (MCES), also called Monte Carlo optimistic policy iteration
(MC-O-PI), is an example: an episode starts from a sampled state--action pair, follows a
greedy policy, and supplies a return used to update the action-value estimate. The policy
is then selected greedily again. A small change in the estimates near a tie can change the
policy, and therefore the target of the next Monte Carlo update.

The convergence question has a long history (Sutton, 1999). Tsitsiklis (2002) established
convergence for discounted optimistic policy iteration with uniform updates, and Chen (2018)
and Liu (2021) studied extensions to stochastic shortest path problems. These works also
distinguish sampling frequencies from effective update rates: component-dependent step sizes
can compensate for nonuniform sampling. A different line of work obtains convergence from
structural assumptions on the MDP, such as an optimal policy feed-forward structure (Wang
et al., 2022). These results concern different combinations of update rules and assumptions;
they do not make initial-visit, first-visit, and every-visit versions interchangeable.

This paper studies initial-visit MCES, in which only the initial state--action pair of each
episode is updated. A common step size suffices when initial action probabilities are equal
within each state and each state accumulates infinite learning weight:
\[
\sigma_k(s,a)=\sigma_k(s)\ge0,\qquad
\sum_{k=0}^{\infty}\alpha_k\sigma_k(s)=\infty\quad(s\in\mathcal S).
\]
The probability $\sigma_k(s)$ is per action. State probabilities may differ, depend on the entire
history, and be zero for arbitrarily long intervals. The update does not divide by the sampling
probability. Infinite visits alone are insufficient as a learning condition: the steps attached to
those visits must have divergent sum.

Theorem~\ref{thm:convergence} covers finite discounted MDPs and proper undiscounted episodic MDPs, with
the return moments and steps specified in Section~\ref{sec:algorithm}. It permits fixed-priority ties or full
inertia, which retains a selected action whenever it remains greedy. Almost surely, $q_k\to q^\star$,
the target $q_{\pi_k}$ changes only finitely often, and all sufficiently late policies are optimal. Policy
labels need not stabilize when optimal actions are tied.

Proposition~12 gives three sampled counterexamples, each changing one algorithmic
choice: the Monte Carlo target, the initial-pair-only update, or equality of initial action
probabilities within a state. Each retains fixed-priority ties, bounded sampled update targets,
a positive sampling floor, and a common nonincreasing Robbins--Monro step sequence. Thus
their failures do not exploit weak state exploration. All three have nonvanishing oscillations
of the optimal-action estimates.

The convergence proof uses the restoring drift of policy-relative gaps. Favorable updates
create positive gaps, a statewise martingale barrier protects them, and a finite ordering of
policy targets forces eventual stabilization. For the counterexamples, robust returns between
separated sets persist on one infinite-horizon noise event.

Section~\ref{sec:preliminaries} specifies the MDP and algorithm, and Section~\ref{sec:results} proves convergence. Section~4
presents the three constructions and Section~5 states their scope. The nonconvergence proofs
and quantitative bounds are in Appendices~A--D; Appendix~E describes the certificates and
finite numerical illustrations.
